\section{随机过程、时间指标与样本路径}
\label{sec:06-Random-Processes/01-Foundations/01}

设想一台传感器，每秒记录一次机房温度。单独看第十秒的读数 \(X_{10}\)——它是一个随机变量，有期望、方差、分布。单独看第二十秒的读数 \(X_{20}\)——它也是一个随机变量。但这一串读数不只是按时间顺序排列的一堆随机变量：第十秒和第二十秒的温度不可能独立。机房的热惯性让相邻读数彼此粘连，前一秒的高温意味着下一秒也大概率偏高。

于是你面临一个新种类的数学对象。它不是单个随机变量，也不是由几个变量组成的固定维度随机向量。它是随时间的推进持续产出的随机数——一个随机过程。

写成序列就是

\[
X_0, X_1, X_2, \ldots
\]

每个 \(X_t\) 是时刻 \(t\) 的随机状态。描述这样一个过程，需要同时回答两个问题：每个时刻单独可能取什么值，以及不同时刻之间如何关联。

\subsection{一个对象，两种固定方式}

随机过程的正式写法是一个二元函数：

\[
X: T \times \Omega \to S,\qquad X(t,\omega) = X_t(\omega).
\]

其中 \(T\) 是时间指标集，\(\Omega\) 是样本空间，\(S\) 是状态空间。这个函数接受一个时间点和一个随机结果，吐出一个状态值。

二元函数给你两种``冻结一个变量''的方式：

\begin{itemize}
\tightlist
\item
  固定时间 \(t\)：\(\omega \mapsto X_t(\omega)\) 是随机变量。对一个固定时刻，概率论的全部工具照常工作——你可以谈 \(X_t\) 的分布、期望、方差。
\item
  固定随机结果 \(\omega\)：\(t \mapsto X_t(\omega)\) 是一条样本路径（或称轨迹）。它是时间的函数，记录了在这一次具体实验中，过程从头到尾走过的轨迹。
\end{itemize}

这两个视角不能混淆。概率分布——无论是有限维联合分布还是路径空间上的测度——是对大量可能路径的总体描述。而现实中一次实验通常只观测到其中一条路径。你看到的是 \(t \mapsto X_t(\omega_0)\) 这一条曲线；你想推断的是产生这条曲线的总体概率律。从一条路径推断总体，需要额外的假设——平稳性、遍历性，后面会专门讲。

\subsection{连续抛硬币：两条路径在同一个过程中}

把上面的抽象框架套到一个最朴素的例子。

连续抛一枚正面概率为 \(p\) 的硬币，设

\[
X_n = \text{前 } n \text{ 次抛掷中的正面总数},\qquad n = 0,1,2,\ldots
\]

固定 \(n=10\)：\(X_{10}\) 是前十次中正面的个数，服从 Binomial\((10,p)\)。这是一个干净的概率分布。

固定一串抛掷结果——比如正反反正正反反正正正——\(n \mapsto X_n\) 是从 \((0,0)\) 出发的一条阶梯路径，每一步要么水平移动一格再向上跳 0 或 1。这条路径不会下降（正面只增不减），斜率时大时小。

同一个过程，两种观察方式得出了不同的画面：固定时间的画面是一个二项分布族，固定结果的画面是一条单调不减的随机阶梯。缺了哪一边，你对过程的感知都是片面的。

这个例子还暗含一个更深的问题：如果只允许写 \(X_t(\omega)\)，还需不需要别的技术条件？通常要求每个固定 \(t\) 的映射 \(X_t: \Omega \to S\) 可测；若要把 \(X\) 当作时间与样本的联合函数来做积分（比如求路径泛函的期望），常常还需要联合可测性。入门计算多从前者出发——先保证每个时刻的随机变量定义合法；连续时间理论才会频繁遭遇联合可测性的技术细节。

\subsection{时间指标与状态空间形成四种基本类型}

时间集合 \(T\) 可以是离散的（\(\{0,1,2,\ldots\}\) 或 \(\mathbb Z\)）或连续的（\([0,\infty)\) 或 \(\mathbb R\)）。状态空间 \(S\) 同样可以是离散的或连续的。四种组合对应四大类模型：

{\def\LTcaptype{none} % do not increment counter
\begin{longtable}{@{}lll@{}}
\toprule\noalign{}
时间 & 状态 & 典型例子 \\
\midrule\noalign{}
\endhead
\bottomrule\noalign{}
\endlastfoot
离散 & 离散 & 每天排队人数、离散时间 Markov 链 \\
离散 & 连续 & 每小时温度记录、ARMA 时间序列 \\
连续 & 离散 & Poisson 计数过程、连续时间 Markov 链 \\
连续 & 连续 & Brownian motion、连续噪声波形 \\
\end{longtable}
}

两个独立的术语很容易产生误解。``连续时间''不是指路径连续——Poisson 过程的时间指标是连续的，但它的样本路径是跳跃的阶梯函数。``连续状态''也不是指时间连续——温度序列可能是逐小时记录的离散时间序列，但温度值本身是连续的量。时间和状态，各管各的。

时间尺度的选择本身就是建模决策。把每秒的事件数聚合成每小时计数，可能抹掉毫秒级的短时依赖结构；把连续变化的物理量以过低的频率采样，可能产生混叠——高频信号伪装成低频信号。不存在唯一正确的时间分辨率；正确的分辨率取决于你关心的物理时间尺度和预测任务的需要。

\subsection{状态应包含预测未来所需的信息}

同一个现实系统可以用不同的状态表示来建模，选择决定了模型能回答什么问题。

以排队系统为例。两种状态表示：

\begin{itemize}
\tightlist
\item
  方案 A：\(X_t = \mathbf 1_{\{\text{队伍非空}\}}\)——队伍是否不为空。模型简单，但无法预测下一次服务完成后的队伍长度。
\item
  方案 B：\(X_t = \text{当前系统中的顾客数}\)——包含了更多信息。给定当前人数，配合到达和服务规则，可以推测下一时刻人数的分布。
\end{itemize}

方案 A 丢掉的信息恰恰是预测未来所需的。队伍从 5 人变成 4 人和从 1 人变成 0 人，在方案 A 中都显示为``非空''，但前者下一时刻大概率仍然非空，后者可能变空。

在 Markov 建模中这件事有一个专门的名称：Markov 性。如果给定当前状态，未来与历史条件独立，就说状态具备 Markov 性。如果未来仍依赖未被状态包含的历史信息——例如当前排队人数相同，但到达率刚刚发生过突变——所选的状态就不满足 Markov 性。常见的修复手段是扩充状态：把``当前位置''扩展为``当前位置 + 速度''，把``队列长度''扩展为``队列长度 + 当前服务已用时间 + 到达率估计''。更丰富的状态向量让过去的信息更多地被当前状态捕获，从而逼近 Markov 性的要求。

\subsection{路径性质与单点分布是不同信息}

两个过程可以在每一个固定时刻拥有完全相同的边缘分布，路径却截然不同。这在直觉上可能有些反直觉——不是每个时刻都一样的分布吗，怎么路径会不一样？

直接构造。设 \(Z \sim \mathcal N(0,1)\) 为一个标准正态随机变量：

\begin{itemize}
\tightlist
\item
  过程 A：对所有 \(t\)，\(X_t = Z\)——整条路径是一条随机选取的水平线。知道任何一个时刻的值，就彻底知道了全部过去和未来的值。
\item
  过程 B：每个整数时刻 \(Y_t\) 独立服从 \(\mathcal N(0,1)\)——路径在各个时刻之间没有任何关联。知道 \(Y_0\) 不告诉你关于 \(Y_1\) 的任何信息。
\end{itemize}

对任何一个固定时刻 \(t\)：

\[
X_t \overset{d}{=} Y_t \sim \mathcal N(0,1),
\]

记号 \(\overset{d}{=}\) 表示两边同分布。只看单个时刻的数字特征——均值、方差、直方图——两个过程完全分不开。但一旦你看两个或更多时刻的联合行为，差别就大到没有边际：过程 A 的路径绝对平坦，过程 B 的路径是锯齿状的完全随机噪声。

因此你不能用一系列边缘分布来代替过程的联合结构。一个过程的信息量远超每个时刻单独分布的简单罗列。

路径的性质还决定了哪些数学操作合法：

\begin{itemize}
\tightlist
\item
  路径是否连续或右连续（关系到随机积分能不能用 Riemann--Stieltjes 型定义）；
\item
  在有限区间内是否只有有限次跳跃（关系到跳跃过程的补偿和变差）；
\item
  路径是否单调（关系到次序统计和极值行为）；
\item
  是否可能在有限时间内爆炸到无穷（关系到停时和过程的全局存在性）；
\item
  是否具有有限变差（关系到路径积分的定义和 Ito 公式是否需要修正项）。
\end{itemize}

这些问题在连续时间过程中会反复回到台前。初学离散时间过程时，你可以暂时把路径性质放在一边——离散时间的路径性质通常比较简单，这些复杂性在连续时间才充分展开。

\subsection{过程模型的第一张清单}

描述一个随机过程时，至少明确以下六项：

\begin{enumerate}
\tightlist
\item
  时间指标 \(T\)——离散还是连续，有无起点或终点；
\item
  状态空间 \(S\)——离散还是连续，是实数、整数还是更抽象的集合；
\item
  初始分布——\(t=0\)（或初始时刻）的状态服从什么分布；
\item
  不同时刻的联合依赖规则——随机变量之间如何关联，这是整个模型最核心也最棘手的部分；
\item
  路径正则性假设——连续、右连续、跳跃次数等；
\item
  参数是否随时间变化——平稳模型（参数恒定）还是非平稳模型（参数是时间的函数）。
\end{enumerate}

前三条通常容易写出。第四条——联合依赖规则——是整个过程的灵魂。下一篇就来回答：如何用有限个时刻的联合分布来系统地描述一个随机过程的依赖结构。
